\section{Tilted nozzle}\label{sec:tilt}

The nozzle can be tilted by an angle $\theta$ (0 to 60°, 0 = normal to the surface) in the plane that contains the path of the nozzle and the normal to the surface. Two things change.

\paragraph{The footprint gets longer, not wider.} A jet of width $\dX$ meets the surface obliquely: its footprint is $\dX/\cos\theta$ long along the path and still $\dX$ wide across it. The flow per unit area under the footprint drops by $\cos\theta$, but a point stays under the jet $1/\cos\theta$ times longer. The swept width is still $\dX$, so the number of shots per unit area in one pass is unchanged:
\begin{equation}
R=\frac{\dot N}{V\,\dX}\quad\text{(independent of }\theta\text{)},\qquad \text{time under the jet } \frac{\dX}{V\cos\theta}.
\end{equation}
The centre of the footprint is also shifted along the path by $d\tan\theta$ ($d$: stand-off distance). This constant shift moves the whole pattern; it does not change the coverage inside the part.

\paragraph{Each dent is smaller.} Only the normal component of the shot velocity, $v_n=v\cos\theta$, indents the surface; the tangential component makes the dent slightly oval and causes some sliding, which the model does not draw. The dent area is multiplied by a factor $f(\theta)$, chosen with the parameter \ui{Dent area at angle θ}:
\begin{equation}
a(\theta)=a_0\,f(\theta),\qquad f(\theta)=\cos\theta\ \ (\text{dent radius}\propto\sqrt{v_n})\quad\text{or}\quad f(\theta)=\cos^2\theta\ \ (\text{normal kinetic energy}),
\end{equation}
where $a_0$ is the dent area at normal incidence (the parameter \ui{Dent radius r} is the radius at normal impact). Measurements usually fall between the two laws. The rate constant becomes $K(\theta)=K_0\,f(\theta)$, and all the times (98\,\%, full coverage, specified coverage) are multiplied by $1/f(\theta)$:

\begin{center}\small
\begin{tabular}{@{}lcccc@{}}
\toprule
$\theta$ & 15° & 30° & 45° & 60°\\ \midrule
time factor, $f=\cos\theta$ & 1.04 & 1.15 & 1.41 & 2.00\\
time factor, $f=\cos^2\theta$ & 1.07 & 1.33 & 2.00 & 4.00\\ \bottomrule
\end{tabular}
\end{center}

\noindent Not modelled: the oval shape of the dents, the lower peening intensity, the fewer collisions between incoming and rebounding shots (which can partly compensate up to about 20--30°), and the grazing impacts at large angles.

\subsection*{Two nozzles}

The parameter \ui{Number of nozzles} adds a second nozzle in the same plane, aimed at the same point and tilted the other way: nozzle~1 at $\theta_1\ge0$, nozzle~2 at $\theta_2\le0$. Each nozzle throws the flow rate $\dot m$ (the parameter is per nozzle). The two jets act independently: their impacts add up, $\lambda=\lambda_1+\lambda_2$, and since each one follows the rules above,
\begin{equation}
K=K_0\,\bigl[f(\theta_1)+f(\theta_2)\bigr],
\end{equation}
where $K_0$ is the rate constant of one nozzle at normal incidence. Two normal nozzles ($\theta_1=\theta_2=0$) are exactly one nozzle with twice the flow rate; two nozzles at $\pm45°$ give $1.41\,K_0$ (law $\cos\theta$) or $1.00\,K_0$ (law $\cos^2\theta$): with the $\cos^2\theta$ law, two nozzles at $\pm45°$ do the work of one normal nozzle. The coverage does not depend on the signs of the angles, only on $f(\theta_1)+f(\theta_2)$ and on the footprint lengths $\dX/\cos\theta_1$ and $\dX/\cos\theta_2$. What two opposite nozzles do change in practice, and the model does not include, is the symmetry of the treatment: the tangential components of the two jets are opposite, so the asymmetric pile-up and the directional effects of a single oblique jet tend to cancel. Collisions between the two streams near the surface are not modelled either.

\subsection*{Angle \texorpdfstring{$\psi$}{psi} across the path}

The angle $\psi$ (0 to 60°) tilts the plane of the nozzle (or of the two nozzles) about the direction of the path. With $\theta$ measured in that plane, the axis of a nozzle, pointing from the surface to the nozzle, is
\begin{equation}
\mathbf a=(\sin\theta,\ \cos\theta\sin\psi,\ \cos\theta\cos\psi)\quad(x\text{ along the path},\ y\text{ across it},\ z\text{ along the normal}),
\end{equation}
so the angle of incidence $\gamma$ is given by $\cos\gamma=\cos\theta\cos\psi$, and the dent area is multiplied by $f(\gamma)=\cos\gamma$ or $\cos^2\gamma$. A jet of width $\dX$ projected along $\mathbf a$ on the surface has a footprint (exact for a Gaussian jet, used for all the profiles)
\begin{equation}
\text{across the path: } \frac{\dX}{\cos\psi},\qquad \text{along the path: } \dX\sqrt{1+\frac{\tan^2\theta}{\cos^2\psi}} .
\end{equation}
The width across the path does not depend on $\theta$, and the length along it reduces to $\dX/\cos\theta$ when $\psi=0$. The consequences differ between the simulators:
\begin{itemize}[nosep,leftmargin=*]
\item \textbf{one line swept again and again} (Shot Peening Coverage Simulator): the swept width grows to $\dX/\cos\psi$, so the impacts per unit area and per pass drop by $\cos\psi$: $K=K_0\cos\psi\sum_k f(\gamma_k)$;
\item \textbf{raster on a flat part}: the mean number of impacts per raster does not depend on the width of the jet, only $f(\gamma)$ changes it; but the wider jet overlaps more lines ($\dX/(\delta\cos\psi)$), which \emph{smooths the stripes};
\item \textbf{cylinder}: the jet no longer aims at the axis. In the cross-section, the ray at a lateral offset $u$ from the jet axis hits the surface at the arc length $s$ given by $u(s)=R\,[\sin(s/R+\psi)-\sin\psi]$; the flux on the surface becomes $g_s(s)=g(u(s))\cos(s/R+\psi)$ for $|s/R+\psi|<\pi/2$, the share of the shots that hit the cylinder becomes $c=G(R(1-\sin\psi))-G(-R(1+\sin\psi))$, and the footprint around the circumference gets longer and asymmetric. With the curvature ignored, $g_s(s)=g(s\cos\psi)\cos\psi$.
\end{itemize}
New views show the angles: the side view along the path shows $\theta_1$, $\theta_2$, and a front view across the path shows $\psi$ (Shot Peening Coverage Simulator and flat part); on the cylinder a cross-section perpendicular to the axis shows the nozzle tilted by $\psi$, its jet, the treated arc of the circumference and the shots that miss. With two nozzles, both lie in the plane tilted by $\psi$, so they appear as one nozzle in the front view or cross-section. The nozzles keep their orientation relative to the path (on the flat part they turn with the 0° and 90° rasters).
